\section{Лебег--Стилтьесийн хэмжээ}
Лебегийн хэмжээг байгуулахдаа интервалын уртыг $b-a$ гэж авсан. Үүний оронд үл буурах функцийн өсөлтийг "урт" болгон авбал илүү ерөнхий хэмжээнүүд гарна. Энэ бүлэгт $F\colon\R\to\R$ нь үл буурах, баруунаасаа тасралтгүй функц байна.

\begin{dfn}\label{lsouter}
$A\subset\R$ бүрийн хувьд
$$\lambda_F^{\star}(A)=\inf\Bigl\{\sum_{i=1}^{\infty}\bigl(F(b_i)-F(a_i)\bigr): A\subset\bigcup_{i=1}^{\infty}(a_i;b_i]\Bigr\}$$
гэж тодорхойлсон $\lambda_F^{\star}$ функцийг $F$ функцээр үүсгэсэн \textbf{Лебег--Стилтьесийн гадаад хэмжээ}\index{Лебег--Стилтьесийн гадаад хэмжээ} гэнэ.
\end{dfn}
$F(x)=x$ үед хагас задгай интервалуудаар хучсан ч Лебегийн гадаад хэмжээ $\ls$ гарна.

\begin{thm}\label{lsmeas}
$\lambda_F^{\star}$ нь $\R$ дээрх гадаад хэмжээ бөгөөд Борелийн олонлог бүр $\lambda_F^{\star}$-хэмжигдэх. Мөн $a\lt b$ бүрийн хувьд
$$\lambda_F^{\star}\bigl((a;b]\bigr)=F(b)-F(a).$$
\end{thm}
\begin{proof}
\textit{Гадаад хэмжээ.} $\lambda_F^{\star}(\vn)=0$ ба монотон чанар тодорхойлолтоос шууд гарна. $A=\bigcup_kA_k$, $\sum_k\lambda_F^{\star}(A_k)\lt\infty$ байг. $\veps\gt 0$ бүрийн хувьд $A_k$-г хучих, $\sum_i(F(b_{k,i})-F(a_{k,i}))\lt\lambda_F^{\star}(A_k)+\veps 2^{-k}$ байх $(a_{k,i};b_{k,i}]$ интервалууд авбал бүх интервалууд нийлээд $A$-г хучих тул $\lambda_F^{\star}(A)\leq\sum_k\lambda_F^{\star}(A_k)+\veps$.

\textit{Хагас интервалын хэмжээ.} $(a;b]$ өөрөө өөрийгөө хучих тул $\lambda_F^{\star}((a;b])\leq F(b)-F(a)$. Эсрэг тэнцэл бишийг харуулъя. $(a;b]\subset\bigcup_i(a_i;b_i]$ ба $\veps\gt 0$ байг. $F$ баруунаасаа тасралтгүй тул $F(a+\delta)\lt F(a)+\veps$ байх $0\lt\delta\lt b-a$, мөн $F(b_i')\lt F(b_i)+\veps 2^{-i}$ байх $b_i'\gt b_i$ олдоно. Тэгвэл $[a+\delta;b]$ компакт олонлог $(a_i;b_i')$ задгай интервалуудаар хучигдах тул Гейне--Борелийн теоремоор төгсгөлөг тооны $(a_{i_1};b_{i_1}'),\ldots,(a_{i_m};b_{i_m}')$ интервалаар хучигдана. Эдгээрийг $a_{i_1}\lt a+\delta$, $b_{i_j}'\gt a_{i_{j+1}}$, $b_{i_m}'\gt b$ байхаар дугаарлаж (илүүдлийг хаяж) болох тул
$$F(b)-F(a+\delta)\leq\sum_{j=1}^{m}\bigl(F(b_{i_j}')-F(a_{i_j})\bigr)\leq\sum_{i=1}^{\infty}\bigl(F(b_i)-F(a_i)\bigr)+\veps.$$
Иймд $F(b)-F(a)\leq\sum_i(F(b_i)-F(a_i))+2\veps$. Доод зааг аваад $\veps\to 0$ гэвэл $F(b)-F(a)\leq\lambda_F^{\star}((a;b])$.

\textit{Борелийн олонлог хэмжигдэх.} Хэмжигдэх олонлогууд $\sigma$-алгебр үүсгэх (Теорем \ref{ghthb}) ба $(-\infty;c]$ хэлбэрийн олонлогууд $\br$-г төрүүлэх тул $E=(-\infty;c]$ хэмжигдэх гэж харуулахад хангалттай. $T\subset\R$, $\lambda_F^{\star}(T)\lt\infty$, $T\subset\bigcup_i(a_i;b_i]$ байг. $I_i=(a_i;b_i]$ интервалыг $I_i\cap E$ ба $I_i\setminus E$ гэсэн (хоосон байж болох) хагас задгай интервалд хуваахад $F$-ийн өсөлт нэмэгдэнэ:
$$F(b_i)-F(a_i)=\bigl(F(\min\{b_i,c\})-F(a_i)\bigr)^{+}+\bigl(F(b_i)-F(\max\{a_i,c\})\bigr)^{+}.$$
$T\cap E\subset\bigcup_i(I_i\cap E)$, $T\setminus E\subset\bigcup_i(I_i\setminus E)$ тул
$$\lambda_F^{\star}(T\cap E)+\lambda_F^{\star}(T\setminus E)\leq\sum_i\bigl(F(b_i)-F(a_i)\bigr).$$
Доод зааг авбал Каратеодорийн нөхцөл $\lambda_F^{\star}(T\cap E)+\lambda_F^{\star}(T\setminus E)\leq\lambda_F^{\star}(T)$ гарна.
\end{proof}

Каратеодорийн теоремоор $\lambda_F^{\star}$-ийн Борелийн $\sigma$-алгебр дээрх зааглалт хэмжээ болно. Үүнийг $\mu_F$ гэж тэмдэглээд $F$ функцээр үүсгэсэн \textbf{Лебег--Стилтьесийн хэмжээ}\index{Лебег--Стилтьесийн хэмжээ} гэнэ. Тодорхойлолтоос $\mu_F$ нь зааглагдсан олонлог бүр дээр төгсгөлөг.

\begin{exm}
\begin{enumerate}[a)]
\item $F(x)=x$ бол $\mu_F=\lambda$ (Борелийн олонлогууд дээр).
\item $F=\chi_{[0;\infty)}$ бол $\mu_F=\delta_0$: $(a;b]$ нь $0$-ийг агуулбал $F(b)-F(a)=1$, эс бөгөөс $0$.
\item $x_k\in\R$, $p_k\geq 0$, $\sum_kp_k\lt\infty$ ба $F(x)=\sum_{x_k\leq x}p_k$ бол $\mu_F=\sum_kp_k\delta_{x_k}$.
\item $F$ нь Канторын функц бол $\mu_F([0;1])=1$, гэвч $\mu_F([0;1]\setminus C)=0$: Канторын олонлогийн гүйцээлтийн интервал бүр дээр $F$ тогтмол. Иймд $\mu_F$ нь Лебегийн тэг хэмжээтэй олонлог дээр төвлөрсөн боловч цэг бүр дээр тэг утгатай.
\end{enumerate}
\end{exm}

\begin{prp}\label{lspoint}
$a\lt b$ бол
$$\mu_F(\{a\})=F(a)-F(a-),\quad \mu_F((a;b))=F(b-)-F(a),\quad \mu_F([a;b])=F(b)-F(a-).$$
Ялангуяа $F$ тасралтгүй байх зайлшгүй бөгөөд хүрэлцээтэй нөхцөл нь $\mu_F$ цэг бүр дээр тэг утгатай байх явдал.
\end{prp}
\begin{proof}
$\{a\}=\bigcap_n(a-\frac1n;a]$ буурах дарааллын огтлол бөгөөд $\mu_F((a-1;a])\lt\infty$ тул хэмжээний дээрээс тасралтгүй чанараар $\mu_F(\{a\})=\lim_n(F(a)-F(a-\frac1n))=F(a)-F(a-)$. Үлдсэн хоёр тэнцэтгэл $(a;b)=(a;b]\setminus\{b\}$, $[a;b]=\{a\}\cup(a;b]$ гэдгээс гарна.
\end{proof}

\begin{thm}\label{lsuniq}
$\mu$ нь $\br$ дээрх, зааглагдсан олонлог бүр дээр төгсгөлөг хэмжээ ба бүх $a\lt b$-ийн хувьд $\mu((a;b])=F(b)-F(a)$ бол $\mu=\mu_F$.
\end{thm}
\begin{proof}
$A\in\br$, $A\subset\bigcup_i(a_i;b_i]$ бол $\mu(A)\leq\sum_i\mu((a_i;b_i])=\sum_i(F(b_i)-F(a_i))$, иймд $\mu(A)\leq\mu_F(A)$. Одоо $A\subset(a;b]$ гэж үзье. Өмнөх тэнцэл бишийг $(a;b]\setminus A$-д хэрэглэвэл
$$\mu(A)=F(b)-F(a)-\mu((a;b]\setminus A)\geq F(b)-F(a)-\mu_F((a;b]\setminus A)=\mu_F(A).$$
Иймд зааглагдсан Борелийн олонлог дээр $\mu=\mu_F$. Дурын $A\in\br$-ийг $A=\bigsqcup_{k\in\Z}A\cap(k;k+1]$ гэж задлаад тоологдом аддитив чанарыг хэрэглэвэл $\mu(A)=\mu_F(A)$.
\end{proof}
Иймд $\R$ дээрх, зааглагдсан олонлог дээр төгсгөлөг Борелийн хэмжээ бүр ямар нэг $F$-ийн Лебег--Стилтьесийн хэмжээ байна: $x\geq 0$ үед $F(x)=\mu((0;x])$, $x\lt 0$ үед $F(x)=-\mu((x;0])$ гэж авахад хангалттай. $F$ ба $F+c$ нэг ижил хэмжээ үүсгэнэ.

\section{Лебег--Стилтьесийн интеграл}
\begin{dfn}\label{lsint}
$f$ нь Борелээр хэмжигдэх функц бөгөөд $\mu_F$ хэмжээгээр интегралчлагдах бол
$$\int_Af\,dF=\int_Af\,d\mu_F\qquad(A\in\br)$$
гэж тэмдэглээд \textbf{Лебег--Стилтьесийн интеграл}\index{Лебег--Стилтьесийн интеграл} гэнэ.
\end{dfn}
Иймд хэмжээгээр интеграл авах тухай бүх теорем (монотон нийлэлт, Фатугийн лемм, зааглагдсан нийлэлт) Лебег--Стилтьесийн интегралд шууд хэрэглэгдэнэ.

\begin{thm}\label{lsrs}
$f$ нь $[a;b]$ дээр тасралтгүй, $a=x_0\lt x_1\lt\ldots\lt x_n=b$ хуваалтын диаметр $d=\max_i(x_i-x_{i-1})$, $\xi_i\in[x_{i-1};x_i]$ бол
$$\lim_{d\to 0}\sum_{i=1}^{n}f(\xi_i)\bigl(F(x_i)-F(x_{i-1})\bigr)=\int_{(a;b]}f\,dF.$$
Өөрөөр хэлбэл тасралтгүй функцийн Риман--Стилтьесийн интеграл Лебег--Стилтьесийн интегралтай тэнцүү.
\end{thm}
\begin{proof}
$g=\sum_if(\xi_i)\chi_{(x_{i-1};x_i]}$ энгийн функцийн интеграл нь Теорем \ref{lsmeas}-ээр яг дээрх нийлбэртэй тэнцүү. $f$ жигд тасралтгүй тул $\veps\gt 0$ бүрийн хувьд $d$ хангалттай бага үед $(a;b]$ дээр $|f-g|\lt\veps$. Иймд
$$\Bigl|\int_{(a;b]}f\,dF-\sum_if(\xi_i)\bigl(F(x_i)-F(x_{i-1})\bigr)\Bigr|\leq\veps\bigl(F(b)-F(a)\bigr).$$
\end{proof}

\begin{thm}\label{lsac}
$g\geq 0$ нь $\R$-ийн зааглагдсан интервал бүр дээр Лебегээр интегралчлагдах ба $F(x)=F(0)+\int_0^xg\,d\lambda$ бол Борелийн $f\geq 0$ функц бүрийн хувьд
$$\int_{\R}f\,dF=\int_{\R}fg\,d\lambda.$$
\end{thm}
\begin{proof}
$\nu(A)=\int_Ag\,d\lambda$ $(A\in\br)$ нь Беппо Левийн теоремоор хэмжээ бөгөөд зааглагдсан олонлог дээр төгсгөлөг, $\nu((a;b])=\int_a^bg\,d\lambda=F(b)-F(a)$. Теорем \ref{lsuniq}-ээр $\nu=\mu_F$. Иймд $f=\chi_A$ үед тэнцэтгэл биелнэ, шугаман чанараар сөрөг биш энгийн функцэд биелнэ. Ерөнхий $f\geq 0$-ийг сөрөг биш энгийн функцүүдийн өсөх дарааллын хязгаар $\varphi_n\uparrow f$ гэж бичээд хоёр талд Монотон нийлэлтийн теоремыг (Теорем \ref{th:15.1}) хэрэглэнэ.
\end{proof}
Ялангуяа $F$ нь абсолют тасралтгүй бол $\int_{(a;b]}f\,dF=\int_a^bf F'\,d\lambda$.

\begin{exm}
\begin{enumerate}[a)]
\item $F(x)=x+\chi_{[1;\infty)}(x)$ бол $\mu_F=\lambda+\delta_1$ тул
$$\int_{(0;2]}x^2\,dF=\int_0^2x^2\,dx+1=\frac{11}{3}.$$
\item $F$ нь Канторын функц бол $F'=0$ бараг бүх цэгт, гэвч $\int_{[0;1]}1\,dF=1\neq 0=\int_0^1F'\,d\lambda$. Иймд Теорем \ref{lsac}-д $F$ абсолют тасралтгүй байх нөхцөл зайлшгүй. Канторын хэмжээ $x\mapsto 1-x$ буулгалтын хувьд тэгш хэмтэй тул $\int_{[0;1]}x\,dF=\int_{[0;1]}(1-x)\,dF$, улмаар $\int_{[0;1]}x\,dF=\frac12$.
\end{enumerate}
\end{exm}
